\label{sec:conclusion}

Administrative zone co-membership is the most legally defensible community character
signal in the M-track, and arguably in the redistricting literature as a whole.
Its legal defensibility stems not from a geographic or economic inference about
community membership but from a direct encoding of the administrative relationships
that courts and commissions have explicitly recognized as communities of interest:
shared school district, shared fire protection authority, shared water utility,
shared electric utility, shared county subdivision or town.

The zone co-membership formula is simple, transparent, and verifiable.
A litigant challenging a map drawn with zone co-membership weights can replicate the
computation from the same public TIGER/Line and EIA Form 861 data.
There are no proprietary data sources, no subjective assessments of community identity,
and no electoral data that could introduce partisan bias.
The formula treats every zone type identically and every tract pair consistently.

\subsection{Expert Witness Language}

An expert witness testifying about a map drawn with zone co-membership weights can
use the following template:
\begin{quote}
  The district boundaries were placed by an algorithm that minimizes cuts through
  edges connecting census tracts that share administrative service zones.
  Specifically, pairs of tracts sharing school district, fire district, water utility,
  electric utility service territory, and county subdivision boundaries were given
  twice the standard cut cost, making splits of such communities geometrically
  costly to the optimization algorithm.

  The zone co-membership score is computed entirely from publicly available data:
  TIGER/Line shapefiles published by the Census Bureau and Form 861 service territory
  data published by the Energy Information Administration.
  No electoral data, no racial composition data, and no incumbency information enters
  the computation.

  The result is that tracts sharing the same school district are much less likely to
  be placed in different congressional districts than they would be under a standard
  geographic weight scheme.
  [For NC-14: no pair of tracts sharing school district, fire district, water utility,
  and electric utility was split across different districts.]
  This outcome is consistent with the communities of interest doctrine as articulated
  in \textit{Reynolds v.\ Sims} and as codified in [applicable state law].
\end{quote}

\subsection{Relationship to M.0 and M.8}

M.6 is one instantiation of the general community character framework established
in M.0.
The zone co-membership formula is not the same as cosine similarity of character
vectors, but it satisfies the same mathematical properties and plugs into the
same edge weight modifier interface.
When the M.8 composite community index combines all M-track signals, zone co-membership
will be one dimension of the composite, weighted by the plan configuration's
\texttt{zone\_alpha} parameter.
The composite index allows a plan designer to emphasize zone co-membership
(perhaps for a commission that has explicitly listed administrative zone integrity
as its primary community criterion) or to blend it with economic character,
housing character, and commuting shed signals for a more holistic community
definition.

\subsection{Limitations}

The zone co-membership formula has three limitations worth acknowledging:
\begin{enumerate}
  \item \textbf{Centroid assignment.}
    The spatial join assigns each tract to one zone per type based on the tract
    centroid.
    Tracts that straddle zone boundaries will be misassigned, potentially
    underestimating co-membership.
    Future work using areal interpolation at the block level would improve accuracy.
  \item \textbf{Incomplete coverage for fire and water districts.}
    The formula degrades gracefully when fire and water district data is unavailable,
    but the degraded formula (with $|Z| = 3$) distinguishes only between
    co-zone and non-co-zone tracts without the granularity provided by
    full zone coverage.
  \item \textbf{Static zone boundaries.}
    Administrative zone boundaries change over time (school district consolidations,
    utility mergers, municipal annexations).
    The formula uses zone boundaries from the most recent TIGER/Line release,
    which may not reflect changes made between census years.
    This is a known limitation shared by all redistricting algorithms that use
    administrative boundaries.
\end{enumerate}

Despite these limitations, zone co-membership provides a more legally grounded and
more verifiable community of interest signal than any qualitative assessment of
neighborhood identity.
It translates a phrase that courts have long recognized---``communities of interest''---into
a formula that courts can verify, litigants can replicate, and redistricting
algorithms can enforce.
